\documentclass[12pt]{article}
\usepackage[utf8]{inputenc}
\usepackage[T1]{fontenc}
\usepackage{hyperref}
\title{Tamanaku Language Resources for K-12 Curriculum Development}
\author{Compiled by Claude Sonnet 5 (claude-sonnet-5) for Andr\'as Kornai}
\date{}
\begin{document}
\maketitle

\section{Introduction}
Tamanaku (ISO 639-3 tmz) is an extinct Cariban language of Venezuela. Glottolog places it in the Mapoyo--Tamanaku branch of Venezuelan Cariban. This report rests on a single round of web searches, and almost nothing exists for the language.

\section{Primary Educational Resources}
None found. The language has no living speakers, so there are no school materials.

\section{Linguistic and Reference Materials}
The earliest record is a word list published by Filippo Salvatore Gilij in 1780. The search results did not give a link to the list itself. Overview: \url{https://en.wikipedia.org/wiki/Tamanaku_language}

Caution: the Wikipedia article on Cumanagoto lists ``Tamanaku'' as an alternate name for that endangered language. Check this against specialist sources before relying on any resource: \url{https://en.wikipedia.org/wiki/Cumanagoto_language}

\section{Religious and Community Resources}
None found.

\section{Digital and Audio Resources}
None found. The Endangered Languages Project has an entry that describes the language as long extinct: \url{https://endangeredlanguages.com/lang/1649?hl=zh-cn}

\section{Curriculum Development Context}
A heritage or revitalization curriculum would have to be rebuilt from the 18th-century word list, and possibly from related Cariban languages. The sources found here do not support teaching materials, so this dossier is short.
\end{document}